% Options for packages loaded elsewhere
% Options for packages loaded elsewhere
\PassOptionsToPackage{unicode}{hyperref}
\PassOptionsToPackage{hyphens}{url}
%
\documentclass[
  ignorenonframetext,
]{beamer}
\newif\ifbibliography
\usepackage{pgfpages}
\setbeamertemplate{caption}[numbered]
\setbeamertemplate{caption label separator}{: }
\setbeamercolor{caption name}{fg=normal text.fg}
\beamertemplatenavigationsymbolsempty
% remove section numbering
\setbeamertemplate{part page}{
  \centering
  \begin{beamercolorbox}[sep=16pt,center]{part title}
    \usebeamerfont{part title}\insertpart\par
  \end{beamercolorbox}
}
\setbeamertemplate{section page}{
  \centering
  \begin{beamercolorbox}[sep=12pt,center]{section title}
    \usebeamerfont{section title}\insertsection\par
  \end{beamercolorbox}
}
\setbeamertemplate{subsection page}{
  \centering
  \begin{beamercolorbox}[sep=8pt,center]{subsection title}
    \usebeamerfont{subsection title}\insertsubsection\par
  \end{beamercolorbox}
}
% Prevent slide breaks in the middle of a paragraph
\widowpenalties 1 10000
\raggedbottom
\AtBeginPart{
  \frame{\partpage}
}
\AtBeginSection{
  \ifbibliography
  \else
    \frame{\sectionpage}
  \fi
}
\AtBeginSubsection{
  \frame{\subsectionpage}
}
\usepackage{iftex}
\ifPDFTeX
  \usepackage[T1]{fontenc}
  \usepackage[utf8]{inputenc}
  \usepackage{textcomp} % provide euro and other symbols
\else % if luatex or xetex
  \usepackage{unicode-math} % this also loads fontspec
  \defaultfontfeatures{Scale=MatchLowercase}
  \defaultfontfeatures[\rmfamily]{Ligatures=TeX,Scale=1}
\fi
\usepackage{lmodern}

\usetheme[]{metropolis}
\ifPDFTeX\else
  % xetex/luatex font selection
\fi
% Use upquote if available, for straight quotes in verbatim environments
\IfFileExists{upquote.sty}{\usepackage{upquote}}{}
\IfFileExists{microtype.sty}{% use microtype if available
  \usepackage[]{microtype}
  \UseMicrotypeSet[protrusion]{basicmath} % disable protrusion for tt fonts
}{}
\makeatletter
\@ifundefined{KOMAClassName}{% if non-KOMA class
  \IfFileExists{parskip.sty}{%
    \usepackage{parskip}
  }{% else
    \setlength{\parindent}{0pt}
    \setlength{\parskip}{6pt plus 2pt minus 1pt}}
}{% if KOMA class
  \KOMAoptions{parskip=half}}
\makeatother


\usepackage{longtable,booktabs,array}
\usepackage{calc} % for calculating minipage widths
\usepackage{caption}
% Make caption package work with longtable
\makeatletter
\def\fnum@table{\tablename~\thetable}
\makeatother
\usepackage{graphicx}
\makeatletter
\newsavebox\pandoc@box
\newcommand*\pandocbounded[1]{% scales image to fit in text height/width
  \sbox\pandoc@box{#1}%
  \Gscale@div\@tempa{\textheight}{\dimexpr\ht\pandoc@box+\dp\pandoc@box\relax}%
  \Gscale@div\@tempb{\linewidth}{\wd\pandoc@box}%
  \ifdim\@tempb\p@<\@tempa\p@\let\@tempa\@tempb\fi% select the smaller of both
  \ifdim\@tempa\p@<\p@\scalebox{\@tempa}{\usebox\pandoc@box}%
  \else\usebox{\pandoc@box}%
  \fi%
}
% Set default figure placement to htbp
\def\fps@figure{htbp}
\makeatother





\setlength{\emergencystretch}{3em} % prevent overfull lines

\providecommand{\tightlist}{%
  \setlength{\itemsep}{0pt}\setlength{\parskip}{0pt}}



 


\setbeamerfont{title}{size=\large} \usepackage{hyperref} \setbeamertemplate{footline}[frame number]
\makeatletter
\@ifpackageloaded{caption}{}{\usepackage{caption}}
\AtBeginDocument{%
\ifdefined\contentsname
  \renewcommand*\contentsname{Table of contents}
\else
  \newcommand\contentsname{Table of contents}
\fi
\ifdefined\listfigurename
  \renewcommand*\listfigurename{List of Figures}
\else
  \newcommand\listfigurename{List of Figures}
\fi
\ifdefined\listtablename
  \renewcommand*\listtablename{List of Tables}
\else
  \newcommand\listtablename{List of Tables}
\fi
\ifdefined\figurename
  \renewcommand*\figurename{Figure}
\else
  \newcommand\figurename{Figure}
\fi
\ifdefined\tablename
  \renewcommand*\tablename{Table}
\else
  \newcommand\tablename{Table}
\fi
}
\@ifpackageloaded{float}{}{\usepackage{float}}
\floatstyle{ruled}
\@ifundefined{c@chapter}{\newfloat{codelisting}{h}{lop}}{\newfloat{codelisting}{h}{lop}[chapter]}
\floatname{codelisting}{Listing}
\newcommand*\listoflistings{\listof{codelisting}{List of Listings}}
\makeatother
\makeatletter
\makeatother
\makeatletter
\@ifpackageloaded{caption}{}{\usepackage{caption}}
\@ifpackageloaded{subcaption}{}{\usepackage{subcaption}}
\makeatother

\usepackage{bookmark}
\IfFileExists{xurl.sty}{\usepackage{xurl}}{} % add URL line breaks if available
\urlstyle{same}
\hypersetup{
  pdfauthor={Giorgio Arcara},
  hidelinks,
  pdfcreator={LaTeX via pandoc}}


\author{Giorgio Arcara}
\date{}

\begin{document}


\begin{frame}
% TITLE SLIDES
\title{Metodologia della Ricerca Clinica\\ \vspace{1em} \emph{05 - Valutare deficit e danni - Formule}}
\author{Giorgio Arcara,\\ Università di Padova \\ IRCCS San Camillo, Venezia}

\titlegraphic{

\vspace*{7cm}
\includegraphics[scale=0.15]{Figures/LogoSanCamilloIRCCS_Unipd_alpha.png}
\begin{center}
\vspace{-2.2em}
\includegraphics[scale=0.15]{Figures/CC_license_3_0.png}
\end{center}
}
%\date{\today}
\maketitle
\end{frame}

\section{5. Valutare i deficit e danni cognitivi -
Formule}\label{valutare-i-deficit-e-danni-cognitivi---formule}

\begin{frame}{Valutare i deficit e danni cognitivi}
\phantomsection\label{valutare-i-deficit-e-danni-cognitivi}
\begin{enumerate}
\tightlist
\item
  Il concetto di soglie per percentuali o probabilità
\item
  Punti z e probabilità
\item
  Percentili
\item
  Campione vs Popolazione
\item
  test per il caso singolo (t di Crawford)
\item
  Metodi di regressione per correzione punteggi
\item
  Metodo dei Punteggi equivalenti
\end{enumerate}
\end{frame}

\begin{frame}{Il cut-off clinico e i valori soglia}
\phantomsection\label{il-cut-off-clinico-e-i-valori-soglia}
Il concetto di cut-off clinico (o in generale di valori soglia o valori
che delimitano certe percentuali) sono fondamentali nella statistica
applicata alla neuropsicologia clinica.

L'idea di fondo, nel loro utilizzo, è trovare un valore che aiuti a
classificare la prestazione del paziente o, più in generale, a
determinare quanto è probabile/improbabile.

L'improbabilità è strettamente legata al concetto di deficit o
presumibile danno cognitivo, che di fatto viene operazionalizzato come
una prestazione insolita, improbabile, assumendo che si sia normali.
\end{frame}

\begin{frame}{Il cut-off clinico e i valori soglia}
\phantomsection\label{il-cut-off-clinico-e-i-valori-soglia-1}
L'utilizzo delle probabilità che viene fatto in ambito clinico è però
insidiosa perché è facile sbagliare nel loro utilizzo e nella loro
interpretazione.

Formalizzare in maniera più rigorosa aiuta ad evitare gli effetti di una
terminologia vaga.
\end{frame}

\begin{frame}{Il cut-off clinico}
\phantomsection\label{il-cut-off-clinico}
\begin{center}
\includegraphics[width=0.5\linewidth,height=\textheight,keepaspectratio]{Figures/Healthy_vs_Pathological.png}
\end{center}

\emph{Non sapendo dove sono popolazioni patologiche (deficit o danno),
mettiamo una soglia arbitrariamente che delimita 5\% più basso (area
grigia). Assumendo che deficit e patologia siano da qualche parte a
sinistra.}
\end{frame}

\begin{frame}{I valori utili per la clinica}
\phantomsection\label{i-valori-utili-per-la-clinica}
Definiamo in generale valori utili alla classificazione clinica quei
punteggi x al di sotto del quale che c'è una probabilità nota k che un
punteggio osservato o abbia un valore uguale o inferiore a x (nei test
in cui un punteggio basso indica una performance peggiore).

\[
\small
\text{definiamo}\ x\ \text{in modo tale che}\ P(o<x) = k
\]

es. con il ``cut-off di normalità'', \(k\) = 0.05 e il valore \(x\) è il
valore tale che \(p(O < x ) = 0.05\) (c'è il 5\% di probabilità di
osservare una prestazione inferiore a x, essendo basso, inferiamo un
danno).

La soglia del 5\% è completamente arbitraria.
\end{frame}

\begin{frame}{I percentili}
\phantomsection\label{i-percentili}
Questa definizione è sovrapponibile a quella che si usa in statistica di
``percentili'', in riferimento a una distribuzione.

L'utilizzo della curva normale o di percentili non parametrici
rappresentano diversi metodi per identificare queste percentuali.

L'uso di punti z non è altro che un metodo (parametrico) che semplifica
l'identificazione di questi percentili.
\end{frame}

\begin{frame}{I percentili}
\phantomsection\label{i-percentili-1}
Ricordiamo che l'utilizzo della curva normale (e \textbf{della
trasformazione in punti z}) ha senso in quanto ci aiuta ad identificare
percentili e cioè valori che delimitano porzioni note della
distribuzione.

\begin{center}
\includegraphics[width=0.8\linewidth,height=\textheight,keepaspectratio]{Figures/gaussiana.png}
\end{center}
\end{frame}

\begin{frame}{Trasformazione in punti z}
\phantomsection\label{trasformazione-in-punti-z}
\[
Z = \frac{X - M}{SD}
\] Ricordiamo che questa formula è utile perchè ci aiuta a definire
percentili (parametrici).

Possiamo usarla in due modi:

\begin{itemize}
\item
  calcoliamo valori z soglia, a partire da media e deviazione standard
  (di solito, a quale punteggio grezzo corrisponde il punteggio
  \[z = - 1.64\], che corrisponde al 5\% percentile).
\item
  trasformiamo il punteggio grezzo in punteggio z e da questo
  determiniamo a che percentile corrisponde.
\end{itemize}
\end{frame}

\begin{frame}{Punti z (altro modo di vedere)}
\phantomsection\label{punti-z-altro-modo-di-vedere}
\small

Quando si trasformano dei punteggi in punti z ci sono alcuni punti che
risultano di importanza.

Se i punteggi \textbf{non} sono standardizzati (non trasformo in punti z
prima).

\begin{itemize}
\tightlist
\item
  Media -- 1.64 * SD delimita il 5\% dei punteggi più bassi
\item
  Media - 1.96 * SD delimita il 2.5\% dei punteggi più bassi
\end{itemize}

Se i punteggi sono \textbf{standardizzati} (punti z con Media = 0 e SD =
1):

\begin{itemize}
\tightlist
\item
  1.64 delimita il 5\% dei punteggi più bassi
\item
  1.96 delimita il 2.5\% dei punteggi più bassi
\end{itemize}

A volte si trova ``2 deviazioni standard''. Non è altro che un altro
numero (sempre arbitrario) per delimitare preformance bassa, sempre
arbitraria (ma più facile da ricordare).
\end{frame}

\begin{frame}{I percentili non parametrici}
\phantomsection\label{i-percentili-non-parametrici}
E nel caso di distribuzione non normali?

Si possono utilizzare \textbf{percentili non parametrici}.

(per esempi vedere codice R associato per potenziale impatto di non
normalità, specie con un marcato effetto soffitto.)
\end{frame}

\begin{frame}{La distribuzione di riferimento}
\phantomsection\label{la-distribuzione-di-riferimento}
Questa definizione di percentili può essere però approfondita
concentrandoci su un punto critico su cui in precedenza abbiamo glissato
e cioè che i valori sono \textbf{relativi alla distribuzione di
riferimento}.

Possiamo quindi specificare meglio la formula così \[
\small
\text{definiamo}\ x\ \text{in modo tale che}\ P(o< x∣N )=k
\] \(x\) è il valore per cui ci si attende un probabilità \(k\) che un
punteggio O sia inferiore a \(x\), assumendo che la persona provenga da
una popolazione di persone ``sane'' (cioè senza deficit cognitivo).

Infatti se la persona venisse da altre popolazioni (es. pazienti con
Alzheimer) i valori attesi potrebbero essere diversi.
\end{frame}

\begin{frame}{Analogia con ipotesi nulla in test statistico}
\phantomsection\label{analogia-con-ipotesi-nulla-in-test-statistico}
\small

Da notare che questa formulazione è del tutto analoga a quella del
p-value in realazione dell'ipotesi nulla e all'identificazione dei
valori critici della statistica (es. il t-value critico oltre il quale
considero il risultato significativo). \[
x : P (t_{obs}<t_{crit}∣H_0)=q
\]

\begin{center}
\begin{footnotesize}
Questa formula è semplificata, e non sono considerati i gradi di libertà
\end{footnotesize}
\end{center}

\(x\) è il valore per cui ci si attende un probabilità \(q\) che la
statistica osservata (\(t_{obs}\) in questo caso) sia inferiore al \(t\)
critico, assumendo che \(H_0\) sia vera.

\(H_0\) dipende dal tipo di test, ma nel caso di un confronto fra
campioni indica che i due campioni che sono stati confrontati con il
t-test provengono dalla stessa popolazione.
\end{frame}

\begin{frame}{Esperimento mentale}
\phantomsection\label{esperimento-mentale}
Supponiamo di avere dei dati normativi e di sapere che la soglia che
delimita il punto \(z\) -1.64 (5\%) è 4, oppure che il percentile 5\% è
4.

Supponiamo di osservare una prestazione di 3.

La consideremmo come deficitaria?

Se vi dicessi che ho calcolato questa soglia da un campione con 20
soggetti o con 200 soggetti, ci fideremmo alla stessa maniera della
soglia?

\textbf{Considerare una soglia per inferenze di tipo clinico non può
prescindere dal fare riflessioni sulla rappresentatività del mio
campione rispetto alla popolazione e su quanto i risultati sono
generalizzabili.}
\end{frame}

\begin{frame}{Campione o popolazione}
\phantomsection\label{campione-o-popolazione}
\begin{columns}
\column{0.8\textwidth}

Un aspetto cruciale è che sebbene percentili siano calcolabili per campione o popolazione, nel contesto della valutazione neuropsicologica l’interesse è sempre (impliciamente) alla popolazione

\textbf{Non mi interessa classificare il mio paziente rispetto al mio specifico campione (es. ai miei 10, 100, 1000 soggetti sani), ma come il mio paziente si classificherebbe in generale rispetto alla popolazione di sani.}

\column{0.2\textwidth}
\begin{figure}
\includegraphics[scale=0.05]{Figures/triangle.png}
\end{figure}
\end{columns}
\end{frame}

\begin{frame}{Campione o popolazione}
\phantomsection\label{campione-o-popolazione-1}
\begin{columns}
\column{0.8\textwidth}

Un altro modo di vedere il problema è che la mia classificazione non dovrebbe dipendere dallo specifico campione normativo, ma dovrebbe valere per tutti i possibili campioni normativi.
\vspace{1em}

Questo vale soprattutto se il mio campione normativo ha dimensioni limitate (ed è quindi poco rappresentativo della popopolazione).

\column{0.2\textwidth}
\begin{figure}
\includegraphics[scale=0.05]{Figures/triangle.png}
\end{figure}
\end{columns}
\end{frame}

\begin{frame}{Percentile o probabilità?}
\phantomsection\label{percentile-o-probabilituxe0}
\begin{columns}
\column{0.8\textwidth}

Notare che il percentile non mi dice nulla di probabilità, ma mi dice il valore che delimita una certa percentuale di osservazioni del campione.
\vspace{1em}

Avere una percentuale non significa necessariamente che sto controllando la probabilità di errore, usare percentili o z-scores si approssima a probabilità (in termini di popolazione) solo quando I campioni sono grandi e sono rispettati assunti (Es. Normalità distribuzione della popolazione).

\column{0.2\textwidth}
\begin{figure}
\includegraphics[scale=0.05]{Figures/triangle.png}
\end{figure}
\end{columns}
\end{frame}

\begin{frame}{Campione o popolazione}
\phantomsection\label{campione-o-popolazione-2}
\begin{center}
\includegraphics[width=0.5\linewidth,height=\textheight,keepaspectratio]{Figures/normativi_percentili.png}
\end{center}

La figura mostra tre ipotetici campioni (n=25) estratti dalla stessa
popolazione.

Il numero sopra indica il cut-off vero (5\%) e quello ottenuto nei
campioni.
\end{frame}

\begin{frame}{Campione o popolazione}
\phantomsection\label{campione-o-popolazione-3}
Usare una soglia calcolata sul 5\% del campione (es. tramite punti z) ha
dunque dei limiti: non mi garantisce che delimiti nella popolazione la
soglia del 5.

Ma possiamo stimare un valore della popolazione a partire da ciò che
abbiamo? (e cioè del campione)

Sì, e questa stima (associata ad errori) è alla base dei concetti di
\textbf{statistica inferenziale}. La statistica inferenziale è (proprio
per definizione) quella branca della statistica che ci permette di fare
inferenze sulla popolazione a partire dai dati del campione, associando
a certe percentuali di errori.
\end{frame}

\begin{frame}{Confronti con valori soglia e statistica inferenziale}
\phantomsection\label{confronti-con-valori-soglia-e-statistica-inferenziale}
\begin{columns}
\column{0.8\textwidth}

Una statistica inferenziale è caratterizzata dal poter fare dei precisi enunciati sulla probabilità di errore riferiti alla popolazione da cui il campione da cui è estratta e non al campione stesso.
\vspace{1em}

In senso generale, non utilizzare un test statistico per fare un inferenza, può essere visto come una strategia qualitativa e meno robusta perchè non è nota la probabilità di commettere errori e la generalizzabilità dei risultati al di là dal campione usato per il confronto.

\column{0.2\textwidth}
\begin{figure}
\includegraphics[scale=0.05]{Figures/triangle.png}
\end{figure}
\end{columns}
\end{frame}

\begin{frame}{Confronti con valori soglia e statistica inferenziale}
\phantomsection\label{confronti-con-valori-soglia-e-statistica-inferenziale-1}
\begin{columns}
\column{0.8\textwidth}

Più concretamente, Se calcolo il cut-off usando la soglia -1.64 del punteggio standardizzato:
\vspace{1em}

1) non so quanto la mia classificazione si generalizzerebbe ad altri campioni.

2) non so quale è la probabilità di errore della mia classificazione (assumendo che la persona sia sana), immaginando infinite ripetizioni.

\column{0.2\textwidth}
\begin{figure}
\includegraphics[scale=0.05]{Figures/triangle.png}
\end{figure}
\end{columns}
\end{frame}

\begin{frame}{Confronti con valori soglia e statistica inferenziale}
\phantomsection\label{confronti-con-valori-soglia-e-statistica-inferenziale-2}
Nota statistica:

Usare il confronto con punti z come test statistico è noto come
\textbf{test-z} e può essere un test statistico, ma assume:

\begin{itemize}
\tightlist
\item
  di \textbf{stare confrontando una media campionaria con una
  popolazione di riferimento}
\item
  di conoscere la \textbf{varianza della popolazione} (non del
  campione).
\end{itemize}

Con varianza non nota si dovrebbe usare il t-test (con la distribuzione
t che si approssima a normale per campioni grandi).
\end{frame}

\begin{frame}{Confronti con valori soglia e statistica inferenziale}
\phantomsection\label{confronti-con-valori-soglia-e-statistica-inferenziale-3}
Usare il confronto con punti z come test statistico implica dunque due
errori:

\begin{enumerate}
[1)]
\tightlist
\item
  sto trattando il mio \emph{singolo individuo come una media
  campionaria}
\item
  sto trattando i dati normativi (che sono un campione) come se fossero
  una popolazione.
\end{enumerate}

Al di là di correttezza formale e teorica, questo implica possibili
differenze nella nostra possibilità di errore di classificare, che in
assenza di statistica appropriata, non è nota.
\end{frame}

\begin{frame}{Nota confortante}
\phantomsection\label{nota-confortante}
In realtà è da riconoscere che con campioni sufficientemente ampi
(\textgreater100) e se la distribuzione è realmente normale e con
deviazione standard non troppo elevate, gli errori non sono drammatici.

(Questa è una considerazione da prendere \emph{cum grano salis}, a
partire da simulazioni dello script 10 associato a queste slides)
\end{frame}

\begin{frame}{t-test di Crawford e Howell}
\phantomsection\label{t-test-di-crawford-e-howell}
Il t-test di Crawford \& Howell (1998)\footnote<.->[frame]{in realtà il
  test è di Sokhal e Rohlf (1995). Crawford \& Howell hanno avuto il
  merito di diffonderlo per I single-case in contesto di
  neuropsicologia.} è una modificazione del t-test che permette di
confrontare un caso singolo rispetto con una media di un campione.

Il t-test di Crawford è una modificazione del t-test che permette di
confrontare un caso singolo rispetto ad un campione, ma con inferenza
rispetto a una popolazione, con errore noto.

\underline{Il test funziona anche con campioni piccoli (~10 soggetti).}
\end{frame}

\begin{frame}{t-test di Crawford \& Howell}
\phantomsection\label{t-test-di-crawford-howell}
\[
t = \frac{X_i - \bar{X_2}}{s_2\sqrt{\frac{N_2+1}{N_2}}}
\]

\vspace{1em}

Dove \(X_i\) è il punteggio del singolo individuo da testare.
\(\bar{X_2}\) è la media del campione normativo. \(S_2\) è la Deviazione
standard del campione normativo ed \(N_2\) è la numerosità del campione
normativo. Confrontando con distribuzione con \(N_2-1\) gradi di libertà
il risultato restituisce la probabilità (data ipotesi nulla vera) di
ottenere valori di \(t\) uguali o inferiori a quelli osservati (vedi
codice di esempio per dettagli).
\end{frame}

\begin{frame}{Note su t-test di crawford (1/2)}
\phantomsection\label{note-su-t-test-di-crawford-12}
Il t-test di Crawford si comporta in maniera più appropriata dello
z-test o di percentili non parametrici (vedi anche script con
simulazioni).

Il t-test di Crawford è anche più robusto in caso di violazione di
normalità. Questo significa che anche se una delle assunzioni è violata,
tutto sommato il comportamento del test è adeguato.

Il test di Crawford funziona bene (a livello di Errore di 1° tipo) anche
per campioni piccoli suggerendo un potenziale utilizzo anche per
campioni raccolti ad-hoc (ma attenzione ad Affidabilità e Validità,
dovrebbero essere disponibili).
\end{frame}

\begin{frame}{Note su t-test di crawford (2/2)}
\phantomsection\label{note-su-t-test-di-crawford-22}
\small

Il t-test di Crawford mantiene un errore di primo tipo costante al 5\%
anche per campioni molto piccoli (n=5).

\textbf{Questo non vuol dire che bastano campioni piccoli, in realtà è
sempre meglio avere campioni numerosi, perchè porteranno cut-off più
reali a quelli della popolazione}

(Si veda simulazione in script 11\_C\&H\_test)

La cosa importante è che il t-test di crawford anche con campioni
piccoli mantiene un errore di 1° tipo a livello del 5\% (concretamente,
per dare un esempio se una persona viene dalla popolazione dei sani,
otterremo un valore sotto ``cut-off'' il 5\% delle volte, proprio come è
atteso. Nel caso del z-score la percentuale potrebbe essere 10\% o più).
\end{frame}

\begin{frame}{Confronto con i dati normativi}
\phantomsection\label{confronto-con-i-dati-normativi}
Il test di Crawford, così come utilizzo di z-score e utilizzo di
percentili, necessita però identificare un campione di riferimento.

I dati dovrebbero essere di soggetti sani (quindi normali) simili alla
persona testa. Ma come definire questo?

\textbf{Spesso per un confronto normativo, ci si riferisce a persone con
stesso sesso e con età, scolarità simili.}

La scelta di queste variabili è pragmatica: sono variabili che sappiamo
influenzare performance cognitiva e sono facili da raccogliere (al
momento di creazione del test).
\end{frame}

\begin{frame}{Confronto con i dati normativi}
\phantomsection\label{confronto-con-i-dati-normativi-1}
\small

Finora abbiamo fatto riferimento a formule che confrontano \textbf{con
un singolo gruppo di dati}, denominato ``campione normativo'' o ``dati
normativi'' che ci permette di ottenere una distribuzione di dati.

\begin{center}
\includegraphics[width=0.4\linewidth,height=\textheight,keepaspectratio]{Figures/normative_data.png}
\end{center}

Questo gruppo sarà fatto di individui con una certa età, una certa
scolarità, etc. e abbiamo sempre sottolineato che I dati normativi
dovrebbero essere simili all'individuo da valutare.

Ma come facciamo a definire se il campione è adeguatamente simile?
\end{frame}

\begin{frame}{Confronto con i dati normativi}
\phantomsection\label{confronto-con-i-dati-normativi-2}
\small

Come facciamo a scegliere un campione di riferimento come simile?

Es. se il mio paziente ha 68 anni e 5 anni di scolarità, quale è un
gruppo adeguato?

Esempi:

\begin{itemize}
\tightlist
\item
  persone da 65 a 75 anni con scolarità tra 3 e 5 anni.
\item
  persone con esattamente 68 anni e esattamente 5 anni di scolarità.
\end{itemize}

In genere i test riportano tabelle predefinite e, a livello di raccolta
di dati è possibile che alcune combinazioni siano poco rappresentate. Al
di là di questo, spesso non è possibile avere accesso agli interi dati e
``comporre'' il dato normativo di riferimento, ma occorre basarsi sulle
tabelle disponibili.
\end{frame}

\begin{frame}{Appropriatezza del campione di riferimento}
\phantomsection\label{appropriatezza-del-campione-di-riferimento}
\small

È un aspetto metodologico importante e di impatto. Se nel mio campione
di riferimento sto usando persone che non sono simili alla persona di
cui sto osservando la performance, le mie inferenze (su deficit o su
danno) sarebbero distorte.

Es. supponiamo che il mio test fornisce un cut-off che si riferisce
all'intervallo 70-80 anni di età.

Supponiamo debba testare un paziente con 70 anni. Vuol dire che nella
distribuzione di dati ci saranno persone di 80 anni di età (la cui
performance potrebbe essere molto più bassa).

\textbf{La stratificazione arbitraria in campioni per le variabili
rilevanti (età, scolarità, sesso), ha sempre problemi sulla composizione
del campione normativo e la sua rappresentatività.}
\end{frame}

\begin{frame}{Metodi di regressione}
\phantomsection\label{metodi-di-regressione}
In alternativa ai metodi che usano un semplice campione di riferimento,
esistono metodi che invece usano metodi basati su regressione.

Il razionale di questi metodi è il seguente:

\begin{enumerate}
[1)]
\tightlist
\item
  A partire da alcune variabili posso prevedere la performance attesa
  del partecipante date le varaibili osservate, secondo un modello di
  regressione lineare
\item
  Una volta che ho questo modello di regressione posso stabilire se la
  performance osservata è sufficientemente improbabile dati i valori
  osservati.
\end{enumerate}
\end{frame}

\begin{frame}{Metodi di regressione}
\phantomsection\label{metodi-di-regressione-1}
Una potenziale formalizzazione dei cut-off ottenuti tramite I metodi di
regressione è la seguente.

Il cut-off \(x\) è il punteggio tale che il paziente i-esimo (con
specifici valori per le variabili aggiuntive V1, V2, \ldots, Vn), tale
che la probabilità di osservare quella performance è \(k\) (con \(k\)
critico normalmente 0.05)

\[
\small
\text{definiamo}\ x\ \text{in modo tale che}\  P(o_i< x∣ V1_i , V2_i , ... , Vn_i , N )=k
\]

Dove \(_i\) indica l'indice per una specifica combinazione di punteggio
osservato data una specifica combinazione di valori di variabili
\((V1, ... V_n)\) che possono influenzare performance Classicamente i
modelli di regressione includono età, scolarità e sesso come variabili
predittori della performance
\end{frame}

\begin{frame}{Metodi di regressione: variabili da considerare}
\phantomsection\label{metodi-di-regressione-variabili-da-considerare}
Quali variabili da considerare nella regressione?

Esistono molti modi per selezionare le variabili da considerare e quali
includere. Approfondire questo è oltre gli scopi di questo corso.

I metodi di regressione più raffinati includono anche metodi per tenere
conto di eventuali ``non-linearità'' (relazioni descrivibili tra
funzioni non lineari, es. Esponenziale, iperbole, e la variabile
osservata).

\begin{center}
\includegraphics[width=0.6\linewidth,height=\textheight,keepaspectratio]{Figures/metodi_regressione.png}
\end{center}
\end{frame}

\begin{frame}{Metodi di regressione: variabili da considerare}
\phantomsection\label{metodi-di-regressione-variabili-da-considerare-1}
\begin{center}
\includegraphics[width=1\linewidth,height=\textheight,keepaspectratio]{Figures/tele_gems.png}
\end{center}

\begin{center}
  \footnotesize
  Esempio da Montemurro et al., 2023
\end{center}
\end{frame}

\begin{frame}{Metodi di regressione}
\phantomsection\label{metodi-di-regressione-2}
Nelle slides che seguono viene mostrato un popolare metodo basato su
regressioni che è alla base della maggior parte dei test
neuropsicologici sviluppati in Italiano, il metodo dei Punteggi
Equivalenti.

Questo metodo contiene due parti principali:

\begin{itemize}
\tightlist
\item
  una parte per ottenere punteggi ``corretti'' tramite regressioni.
\item
  una parte per ottenere delle soglie che tengano conto della differenza
  campione/popolazione
\end{itemize}
\end{frame}

\begin{frame}{I Punteggi Equivalenti}
\phantomsection\label{i-punteggi-equivalenti}
Introdotti da Erminio Capitani nella raccolta di taratura di test di
Spinnler e Tognoni (1987)

Tre aspetti chiave:

\begin{itemize}
\tightlist
\item
  Correggono per età e scolarità tramite regressioni.
\item
  Tengono conto di differenza campione/popolazione
\item
  Permettono il confronto tra diversi test
\end{itemize}
\end{frame}

\begin{frame}{I Punteggi Equivalenti}
\phantomsection\label{i-punteggi-equivalenti-1}
Il calcolo dei Punteggi equivalenti è suddiviso in vari step:

\begin{enumerate}
[1)]
\tightlist
\item
  Calcolo di punteggi aggiustati, a cui è rimosso l'effetto di età e
  scolarità.
\item
  Calcolo limite di tollerenza (PE = 0).
\item
  Calcolo dei rimanenti PE.
\end{enumerate}
\end{frame}

\begin{frame}{I Punteggi Equivalenti}
\phantomsection\label{i-punteggi-equivalenti-2}
\textbf{1) Calcolo di punteggi aggiustati, a cui è rimosso l'effetto di
età e scolarità (e sesso).}

\begin{columns}
\column{0.5\textwidth}

La distribuzione dei punteggi grezzi include persone di diversa età.

\column{0.5\textwidth}
\begin{figure}
\includegraphics[scale=0.25]{Figures/raw_scores_distr_normative_data.png}
\end{figure}
\end{columns}
\end{frame}

\begin{frame}{I Punteggi Equivalenti}
\phantomsection\label{i-punteggi-equivalenti-3}
\textbf{1) Calcolo di punteggi aggiustati, a cui è rimosso l'effetto di
età e scolarità (e sesso).}

\begin{columns}
\column{0.5\textwidth}

Effetto dell'età

\column{0.5\textwidth}
\begin{figure}
\includegraphics[scale=0.25]{Figures/age_effect.png}
\end{figure}
\end{columns}
\end{frame}

\begin{frame}{I Punteggi Equivalenti}
\phantomsection\label{i-punteggi-equivalenti-4}
\textbf{1) Calcolo di punteggi aggiustati, a cui è rimosso l'effetto di
età e scolarità (e sesso).}

\begin{columns}
\column{0.5\textwidth}

Effetto dell'età e retta di regressione

\column{0.5\textwidth}
\begin{figure}
\includegraphics[scale=0.25]{Figures/age_effect_regr_line.png}
\end{figure}
\end{columns}
\end{frame}

\begin{frame}{I Punteggi Equivalenti}
\phantomsection\label{i-punteggi-equivalenti-5}
\textbf{1) Calcolo di punteggi aggiustati, a cui è rimosso l'effetto di
età e scolarità (e sesso).}

\begin{columns}
\column{0.5\textwidth}

Differenza tra punteggio previsto e punteggio osservato

\column{0.5\textwidth}
\begin{figure}
\includegraphics[scale=0.25]{Figures/age_effect_regr_line_residual.png}
\end{figure}
\end{columns}
\end{frame}

\begin{frame}{I Punteggi Equivalenti}
\phantomsection\label{i-punteggi-equivalenti-6}
\textbf{1) Calcolo di punteggi aggiustati, a cui è rimosso l'effetto di
età e scolarità (e sesso).}

\begin{columns}
\column{0.5\textwidth}

Calcolo residui di tutti i soggetti dei dati normativi

\column{0.5\textwidth}
\begin{figure}
\includegraphics[scale=0.25]{Figures/age_effect_regr_line_allresidual.png}
\end{figure}
\end{columns}
\end{frame}

\begin{frame}{I Punteggi Equivalenti}
\phantomsection\label{i-punteggi-equivalenti-7}
\textbf{1) Calcolo di punteggi aggiustati, a cui è rimosso l'effetto di
età e scolarità (e sesso).}

\begin{columns}
\column{0.5\textwidth}

I residui potrebbero essere considerati già punteggi «aggiustati» perché sono ciò che rimane del punteggio tolta l’influenza dei predittori (in questo caso ripulito dell’effetto di età)

\column{0.5\textwidth}
\begin{figure}
\includegraphics[scale=0.25]{Figures/residual_distr_normative_data.png}
\end{figure}
\end{columns}
\end{frame}

\begin{frame}{I Punteggi Equivalenti}
\phantomsection\label{i-punteggi-equivalenti-8}
\textbf{1) Calcolo di punteggi aggiustati, a cui è rimosso l'effetto di
età e scolarità (e sesso).}

\begin{columns}
\column{0.5\textwidth}

Nel metodo dei PE, i punteggi aggiustati sono riportati in una scala simile a quella originale (tecnicamente, aggiungendo l’intercetta al residuo)

\column{0.5\textwidth}
\begin{figure}
\includegraphics[scale=0.25]{Figures/adj_scores_distr_normative_data.png}
\end{figure}
\end{columns}
\end{frame}

\begin{frame}{I Punteggi Equivalenti}
\phantomsection\label{i-punteggi-equivalenti-9}
\textbf{1) Calcolo di punteggi aggiustati, a cui è rimosso l'effetto di
età e scolarità (e sesso).}

\begin{columns}
\column{0.5\textwidth}

Si calcola il PE = 0 come \textbf{limite di tolleranza interno unidirezionale non parametrico (5\%)}

Indica il punteggio che non verrà ottenuto da più del 5\% della popolazione. (quindi il cut-off, a livello di popolazione)

\column{0.5\textwidth}
\begin{figure}
\includegraphics[scale=0.25]{Figures/limiti_tolleranza.png}
\end{figure}
\end{columns}
\end{frame}

\begin{frame}{I Punteggi Equivalenti}
\phantomsection\label{i-punteggi-equivalenti-10}
\textbf{2) Calcolo limite di tolleranza (PE = 0)}

\begin{columns}
\column{0.5\textwidth}

In riferimento all’ipotetica distribuzione dell’abilità sottostante, si trova anche il PE=4 e poi si divide lo spazio tra PE = 0 e PE = 4 in intervalli equispaziati. (equispaziati rispetto all’abilità sottostante)

\column{0.5\textwidth}
\begin{figure}
\includegraphics[scale=0.25]{Figures/PE_distr.png}
\end{figure}
\end{columns}
\end{frame}

\begin{frame}{I Punteggi Equivalenti}
\phantomsection\label{i-punteggi-equivalenti-11}
\textbf{2) Calcolo limite di tolleranza (PE = 0)}

\begin{columns}
\column{0.5\textwidth}

1) Calcoliamo il \textbf{punteggio aggiustato} tramite confronto con tabella o applicazione della formula di regressione

2) Confronto con tabella di corrispondenza punteggio aggiustato – Punteggio equivalente

Non si correggono risultati che ottengono il punteggio massimo o minimo.

\textbf{Se punteggio corretto è < soglia PE = 0, allora prestazione deficitaria.}

\column{0.5\textwidth}
\begin{figure}
\includegraphics[scale=0.25]{Figures/adj_score.png}
\end{figure}
\end{columns}
\end{frame}

\begin{frame}{I Punteggi Equivalenti}
\phantomsection\label{i-punteggi-equivalenti-12}
Proviamo insieme a fare un ipotetico utilizzo di punteggi equivalenti.

(vedere da articolo)
\end{frame}

\begin{frame}{Considerazioni su metodi di regressione}
\phantomsection\label{considerazioni-su-metodi-di-regressione}
I metodi di regressioni possono sembrare una soluzione perfetta al
problema sollevato per i test più semplici di confronto con dati
normativi (es. T-test di Crawford). In particolare \textbf{evitano
gradini arbitrari} e permettono una stima più corretta del paziente.

Non sono però privi di problemi. In particolare una stima non corretta
del modello di regressione che porta ai punteggi finali può portare a
stime distorte.
\end{frame}

\begin{frame}{Considerazioni su metodi di regressione}
\phantomsection\label{considerazioni-su-metodi-di-regressione-1}
In generale ci aspettiamo effetti di distorsione maggioramente per
valori estremi nelle variabili predittori, esempio scolarità o età molto
alte o molto basse (sono i punti in cui eventuali imprecisioni del
modello darebbero effetti più accentuati).

I modelli di regressione possono inoltre dare l'illusione di avere delle
buone predizioni per alcuni punti (nello spazio dei valori possibili)
per cui in realtà potrebbero esserci poche osservazioni e stime
imprecise.
\end{frame}

\begin{frame}{Limiti dei metodi di regressione (1/2)}
\phantomsection\label{limiti-dei-metodi-di-regressione-12}
I modelli di regressione possono inoltre dare l'illusione di avere delle
buone predizioni per alcuni punti (nello spazio dei valori) per cui in
realtà potrebbero esserci poche osservazioni e stime imprecise.

Ad esempio, supponiamo di avere dati normativi di 1000 soggetti con un
metodo di regressione. Supponiamo però di dovere testare un soggetto
particolare (es. Donna con 80 anni e 23 anni di scolarità).
\end{frame}

\begin{frame}{Limiti dei metodi di regressione (1/2)}
\phantomsection\label{limiti-dei-metodi-di-regressione-12-1}
È possibile che in realtà nel mio campione ci siano pochissimi (o
addirittura nessuna, in tal caso parliamo di \emph{estrapolazione})
persona che rispecchi questa combinazione di valori.

Sarebbe più opportuno avere un piccolo campione (es. anche 15 persone)
che però siano veramente rappresentative rispetto all'individuo che
voglio valutare e un test appropriato (es. T-test di Crawford).
\end{frame}

\begin{frame}{Considerazioni su errore di I e II tipo}
\phantomsection\label{considerazioni-su-errore-di-i-e-ii-tipo}
\begin{columns}
\column{0.8\textwidth}

In generale abbiamo visto come usare metodi che tengono conto di differenza tra campione e popolazione sia meglio.

Il metodo t di Crawford (e metodi di regressioni di Crawford) mostrano anche che sono in grado di mantenere l’errore di I tipo (cioè avere un soggetto sotto cut-off, quando è sano), a livello atteso (cioè 0.05 o 5%) quando ci sono un numero elevato (in realtà infinito) di ripetizioni.

\column{0.2\textwidth}
\begin{figure}
\includegraphics[scale=0.05]{Figures/triangle.png}
\end{figure}
\end{columns}
\end{frame}

\begin{frame}{Considerazioni su errore di I e II tipo}
\phantomsection\label{considerazioni-su-errore-di-i-e-ii-tipo-1}
\begin{columns}
\column{0.8\textwidth}

Esiste anche un altro tipo di errore rilevante che è quello di II tipo, cioè non riuscire ad identificare correttamente un soggetto patologico. Per i test utilizzati per il deficit/danno, questo è difficile da calcolare, perchè noi non sappiamo quale è la performance attesa per chi ha deficit/danno (vedere figura in slides iniziali). 

Per calcolare quindi errore di primo tipo, si studiano in genere diversi possibili scenari \textbf{ipotetici} con diverse gravità di deficit. In genere però questo \underline{non è mai fatto} per test neuropsicologici. E ci si concentra solo sul cut-off e sull'errore di I tipo.

\column{0.2\textwidth}
\begin{figure}
\includegraphics[scale=0.05]{Figures/triangle.png}
\end{figure}
\end{columns}
\end{frame}

\begin{frame}{Considerazioni su errore di I e II tipo}
\phantomsection\label{considerazioni-su-errore-di-i-e-ii-tipo-2}
\begin{columns}
\column{0.8\textwidth}

Errore di I e II tipo inoltre si applicano non ad uno specifico campione, ma immaginando \textbf{infiniti} campioni. 

Si veda simulazione per vedere quanto è più variabile la stima del cut-off con campioni piccoli.

Viste questi aspetti sarebbe sempre desiderabile (laddove possibile) avere campioni ampi.

\column{0.2\textwidth}
\begin{figure}
\includegraphics[scale=0.05]{Figures/triangle.png}
\end{figure}
\end{columns}
\end{frame}

\begin{frame}{Considerazione}
\phantomsection\label{considerazione}
\pause
\Large
\center

Che fine hanno fatto affidabilità e validità?
\end{frame}

\begin{frame}{Affidabilità, Validità e cut-off}
\phantomsection\label{affidabilituxe0-validituxe0-e-cut-off}
Come anticipato, per convenzione spesso nel calcolo dei valori soglia
affidabilità e validità non sono considerate. \pause

Questo è un problema perchè test con affidabilità diversa sono trattati
allo stesso modo (ho work-in-progress su questo).

Affidabilità, Validità e cut-off/trasformazioni in percentili sono
dunque combinate solo tramite interpretazione dei risultati e non
tramite formule matematiche che tengono in considerazione entrambe.
\end{frame}

\begin{frame}{Riassunto finale (1/2)}
\phantomsection\label{riassunto-finale-12}
Siamo partiti dal problema di definire deficit e danno.

Il metodo più diffuso è basato sull'improbabilità di osservare quella
prestazione se si è sani (ragionamento contorto, ma necessario, vedi
slides).

Spesso c'è necessità di una soglia, ed essa è arbitrariamente messa al
5\% delle prestazioni più basse. Da notare che entrambe le scelte
ricordano il rifiuto dell'ipotesi nulla che c'è per i test statistici.
\end{frame}

\begin{frame}{Riassunto finale (1/2)}
\phantomsection\label{riassunto-finale-12-1}
\begin{center}
PROBLEMA: cosa si intende per 5\%? 
\end{center}

Non avrebbe senso riferirsi alla distribuzione del campione, ma sarebbe
meglio riferirsi alla popolazione.

In generale esistono vari metodi per calcolare soglie e probabilità, e
sarebbe meglio:

\begin{enumerate}
[a)]
\item
  usare quelli che tengono effettivamente conto di distinzione
  campione/popolazione
\item
  avere campioni numerosi (almeno n \textgreater{} 100).
\end{enumerate}

Anche usando metodi non ottimali, avere campioni ampi migliora.
\end{frame}

\begin{frame}{Riassunto finale (2/2)}
\phantomsection\label{riassunto-finale-22}
I metodi che si basano su campioni hanno però problemi dei ``gradini''
che si vengono a creare per ogni gruppo in cui abbiamo trasformato una
variabile continua (es. Età) in gruppi separati.

I metodi che si basano su regressione risolvono il problema dei
``gradini'', ma attenzione a rappresentatività nei metodi con
regressione e attenzione a valori estremi nei predittori (es. Scolarità,
età). Sono I valori in cui la stima potrebbe essere più imprecisa.
\end{frame}




\end{document}
